\documentclass[12pt]{article}
\usepackage[margin=1in]{geometry}
\usepackage{parskip}
\usepackage{microtype}

\begin{document}
\pagestyle{empty}

\today

\bigskip

Dear Editor,

We are pleased to submit our manuscript titled ``Selective Prediction and the Persistence Illusion: A Diagnostic Decomposition of VIX Regime Classification'' for consideration for publication in the \textit{Journal of Risk and Financial Management}.

This paper addresses a methodological issue of direct relevance to practitioners and researchers in financial risk management: the reliability of machine learning-based regime classification systems. We develop a six-component diagnostic protocol and apply it to VIX regime prediction across 5,000 trading days (July 2006--May 2026), demonstrating that confidence-based selective prediction achieves apparent accuracy of 90--93\% not through genuine predictive skill, but through a mechanism we term the ``persistence illusion''---the concentration of predictions on autocorrelation-stable days.

Our key findings are:

\begin{enumerate}
  \item A three-feature HAR-style logistic regression matches or exceeds a 36-feature Random Forest at four of five prediction horizons under bootstrap testing, meaning the additional 33 features contribute negligibly.

  \item A McNemar test on jointly covered days finds $b = c = 0$ at $N \geq 10$: the Random Forest and a naive persistence rule make identical predictions on every shared day, confirming the accuracy gap is attributable entirely to coverage composition, not predictive skill.

  \item The confidence filter achieves 0\% accuracy on calm-to-high volatility transitions at all horizons. Explicit cost-weighting up to $20\times$ does not improve this, establishing an informational ceiling rather than an objective-function deficiency.

  \item The full diagnostic replicates on S\&P 500 trend regimes, establishing the persistence illusion as a general structural property of selective prediction on serially correlated labels.
\end{enumerate}

We believe this work is well-suited to JRFM given its focus on practical risk management implications and its accessible, modular methodology. The paper contributes both a negative empirical result and a positive methodological contribution: a reusable six-component diagnostic protocol applicable to any selective prediction system on autocorrelated outcomes.

All data are publicly available (Yahoo Finance and FRED). Code is available upon request.

We confirm that neither the manuscript nor any parts of its content are currently under consideration for publication with or published in another journal.

All authors have approved the manuscript and agree with its submission to JRFM.

We thank you for your time and consideration.

\bigskip

Sincerely,

\bigskip\bigskip

Akshat Gupta\\
Texas Academy of Mathematics and Science, University of North Texas\\
akshatguptaakshatgupta@gmail.com\\[6pt]
Dr.\ Jianguo Liu (Corresponding Author)\\
Department of Mathematics, University of North Texas\\
jianguo.liu@unt.edu

\end{document}
